The Value Proposition explains how the sponsors will benefit from your work, and why they should invest funding, time, and expertise in supporting your team. Here, you are essentially making a case for the project. There are many ways in which value can be returned to your stakeholders (industrial sponsors, instructors, the university, etc.), list any that may help you convince them to "buy in".

The 3D Lab Manager gives Lab 208 in the Engineering Research Building a single system that shows what equipment exists,
where it sits on a 3D scan of the lab, and whether it is available right now. 
A touchscreen kiosk in the lab puts this information in front of students when they need it. 
A previous team attempted this system but never completed it.

The project returns value to its stakeholders in many ways:

\begin{itemize}
    \item \textbf{Time returned to the lab manager.}
    Students can find items on their own instead of contacting
    Steven McDermott, freeing him for his other responsibilities.

    \item \textbf{Protection of high value equipment.}
    Wireless tracking shows where expensive tools are, reducing
    loss and protecting the department's investment.

    \item \textbf{Recognition for student work.}
    A project showcase highlights creative work built in Lab 208 for prospective students to see
    and credits the students who made it.

    \item \textbf{Promotion of UTA.}
    Public 3D scans of UTA labs let prospective students,
    sponsors, and visitors explore the university's facilities.

    \item \textbf{Long term value.}
    The system runs on university servers and keeps serving
    future senior design teams at no licensing cost.
\end{itemize}